\documentclass{article}
\usepackage[T1]{fontenc}
\usepackage{lmodern}
\usepackage{graphicx}
\usepackage{amsmath,amssymb}
\usepackage[a4paper,margin=2cm]{geometry}
\usepackage[absolute,overlay]{textpos}
\setlength{\TPHorizModule}{1mm}
\setlength{\TPVertModule}{1mm}
\usepackage[indonesian]{babel}
\usepackage{hyperref}
\setlength{\emergencystretch}{2em}
% unit: Halaman 47

\begin{document}

% === Halaman 47 (python/native) ===

47

def dekripsi\_aes(cipherteks, kunci):

    kunci = kunci.encode()                        \# Kodekan kunci dari string menjadi byte

    cipherteks = base64.b64decode(cipherteks)     \# Dekodekan kembali cipherteks dari base64 ke byte 

    iv = cipherteks[:16]                          \# Ambil iv pada bagian awal cipherteks (16 byte)

    AEScipher = AES.new(kunci, AES.MODE\_CBC, iv)  \# Defenisikan AES mode CBC 

    plainteks = AEScipher.decrypt(cipherteks[BLOCK\_SIZE:])     \# Lakukan dekripsi cipherteks

    plainteks = unpad(plainteks, BLOCK\_SIZE)      \# Buang padding

    return plainteks  

\# Program utama (pemanggilan fungsi enkripsi dan dekripsi AES)

\# Enkripsi pesan 

print("E N K R I P S I:")

pesan = input("Ketikkan pesan yang akan dienkripsi:")

kunci = input("Ketikkan kunci (harus sepanjang 16 karakter):")

cipherteks = enkripsi\_aes(pesan, kunci)

print("Cipherteks hasil enkripsi: ", cipherteks)

\# Dekripsi pesan

print("D E K R I P S I:")

kunci = input("Ketikkan kunci (harus sepanjang 16 karakter):")

plainteks = dekripsi\_aes(cipherteks, kunci)

\#plainteks = unpad(plainteks, BLOCK\_SIZE)

print("Plainteks hasil dekripsi: ", plainteks[16:])

\end{document}
